\documentclass[10pt,a4paper]{article}
\usepackage[margin=1.5cm]{geometry}
\usepackage{times}
\usepackage{array}
\usepackage{amsmath,amssymb,amsfonts}
\usepackage{booktabs}
\usepackage{enumitem}
\usepackage{hyperref}
\usepackage[table]{xcolor}
\usepackage{titlesec}

\hypersetup{colorlinks=true,linkcolor=black,urlcolor=blue,citecolor=black}
\setlength{\parindent}{0pt}
\setlength{\parskip}{2.5pt}
\setlist[itemize]{leftmargin=1.15em,itemsep=1pt,topsep=2pt}
\setlist[enumerate]{leftmargin=1.25em,itemsep=1pt,topsep=2pt}
\renewcommand{\arraystretch}{1.08}
\titleformat{\section}{\normalsize\bfseries}{\thesection}{0.55em}{}
\titlespacing*{\section}{0pt}{5pt}{2pt}
\titleformat{\subsection}{\normalsize\bfseries}{\thesubsection}{0.45em}{}
\titlespacing*{\subsection}{0pt}{4pt}{1.5pt}

\newcommand{\blank}[1][7em]{\colorbox{yellow!35}{\rule{0pt}{1.05em}\hspace{#1}}}
\newcommand{\eg}{E_{\mathrm{g}}}
\newcommand{\ehull}{E_{\mathrm{hull}}}
\newcommand{\form}{x_{\mathrm{form}}}
\newcommand{\struc}{x_{\mathrm{struc}}}
\newcommand{\unc}{u}
\newcommand{\domain}{d_{\mathrm{dom}}}
\newcommand{\utility}{U}
\newcommand{\gapscore}{\rho_{\mathrm{gap}}}
\newcommand{\appscore}{\rho_{\mathrm{app}}}
\newcommand{\novel}{\rho_{\mathrm{novel}}}

\begin{document}
\vspace*{-34pt}
\begin{center}
    \textbf{CITY UNIVERSITY OF HONG KONG}\\
    \vspace{6pt}
    \underline{\textbf{Application Form for Huawei Seed Grant}}
\end{center}

\section{Name of the applicant, position and home Department, and the date of joining the University}
\begin{itemize}[leftmargin=15pt]
    \item Name of applicant: Chen MA
    \item Position and affiliated department: Assistant Professor, Department of Computer Science
    \item Date of joining CityU: August 25, 2021
    \item Duration of the project: \blank[18em]
\end{itemize}

\section{Title of the Proposed Research}
\begin{itemize}[leftmargin=15pt]
    \item English Project Title: \textbf{AI-Powered Boron Nitride Material Exploration}
    \item Chinese Project Title: \blank[18em]
\end{itemize}

\section{Abstract of the Project}
AI is increasingly used to accelerate materials discovery, but BN-related materials remain difficult to explore because available labels are sparse, heterogeneous, and uneven across bulk, 2D, doped, and carbon/silicon-containing variants. This project aims to build an AI-powered exploration framework for boron nitride materials. The framework will construct a BN-centered data layer, learn formula-only and structure-resolved property models with calibrated uncertainty, rank candidates under application and reliability constraints, and generate prototype-based structure hypotheses for validation handoff. The project will deliver a decision pipeline that identifies which BN families should be investigated next for wide-band-gap, dielectric, and two-dimensional-material applications. Expected outputs include a curated dataset, benchmark report, ranked candidate table, structure package, and manuscript-ready analysis.

\section{Key Objectives}
The central research question is:
\[
  \textit{How can AI transform scattered BN-related materials data into reliable, uncertainty-aware research decisions?}
\]
The project will address this question through four objectives.
\begin{itemize}[leftmargin=*]
  \item \textbf{Build a bounded BN exploration space.} The candidate space will be organized as
  \[
    \mathcal{C}_{\mathrm{BN}}=
    \mathcal{C}_{\mathrm{core}}\cup\mathcal{C}_{\mathrm{ext}}\cup\mathcal{C}_{\mathrm{explore}},
  \]
  where $\mathcal{C}_{\mathrm{core}}$ covers BN, h-BN, c-BN, and doped BN;
  $\mathcal{C}_{\mathrm{ext}}$ covers BCN, BC$_2$N, and SiBN; and
  $\mathcal{C}_{\mathrm{explore}}$ is restricted to charge-neutral B/N/C/Si-centered substitutions that map to known BN-family motifs:
  \[
    \mathcal{C}_{\mathrm{explore}}=
    \{c:q(c)=0,\;N_{\mathrm{atom}}(c)\le 12,\;\pi(c)\in
    \{\mathrm{BN},\mathrm{BCN},\mathrm{BC_2N},\mathrm{SiBN}\}\}.
  \]
  \textit{Output:} a provenance-labelled BN dataset and candidate design space.

  \item \textbf{Develop two-stage BN property models.} Formula-only candidates will first be screened using composition descriptors. Candidates passing the first screen will enter structure hypothesis generation, after which graph or structure-resolved models will be applied:
  \[
    f_{\theta}^{(1)}:\form\rightarrow(\hat{\eg},\hat{\ehull},\hat{\unc}),
    \qquad
    f_{\theta}^{(2)}:(\form,\struc)\rightarrow
    (\hat{\eg},\hat{\ehull},\hat{\epsilon}_r,\hat{\unc}).
  \]
  \textit{Output:} benchmarked formula-only and structure-resolved models with calibrated uncertainty.

  \item \textbf{Rank candidates through an uncertainty-aware decision layer.} Candidate utility will combine target-window property scores, stability, uncertainty, domain distance, novelty, and application relevance:
  \[
    \utility(c)=w_1\gapscore(c)-w_2\hat{\ehull}(c)-w_3\hat{\unc}(c)
    -w_4\domain(c)+w_5\novel(c)+w_6\appscore(c).
  \]
  \textit{Output:} a ranked candidate table with predicted properties, uncertainty, novelty score, application label, and recommended action.

  \item \textbf{Create validation-ready structure hypotheses.} Priority formulas will be converted into structure candidates through prototype reuse, database-derived structure transfer, and controlled substitution before any generative route is considered:
  \[
    c\mapsto\{s_{c,1},\ldots,s_{c,k}\}\mapsto
    \{\mathrm{sanity\ check},\mathrm{relaxation},\mathrm{property\ recheck}\}.
  \]
  \textit{Output:} CIF/POSCAR files, provenance records, predicted properties, uncertainty scores, and validation notes.
\end{itemize}

\section{Research Methodology}
\subsection{Background of Research}
Boron nitride is a strategically important materials family because its polymorphs and derivatives support wide-band-gap optoelectronics, dielectric interfaces, high-quality 2D heterostructures, and chemically robust insulating layers. h-BN and c-BN have been extensively studied, while BCN, BC$_2$N, SiBN, doped BN, and related substituted analogues provide a wider design space for AI-guided exploration. However, BN-related labels are sparse and unevenly distributed across public materials databases. A naive model trained on broad materials data may score well on standard splits but fail on BN-family extrapolation. Therefore, the proposed research treats BN not as a single compound, but as a target materials family requiring family-aware data construction, modelling, ranking, and validation.

\subsection{WP1: BN-Centered Data Construction}
The project will integrate public computational resources including 2DMatPedia, JARVIS, Matbench, matminer, the Materials Project, and C2DB~\cite{zhou2019_2dmatpedia,choudhary2020_jarvis,dunn2020_matbench,ward2018_matminer,jain2013_mp,haastrup2018_c2db}. Each record will be represented as
\[
  r_i=(c_i,s_i,\mathbf{y}_i,p_i),\qquad
  p_i=(\mathrm{source},\mathrm{method},\mathrm{version},\mathrm{provenance}),
\]
where $c_i$ is composition, $s_i$ is optional structure, $\mathbf{y}_i$ is the property vector, and $p_i$ stores data lineage. The BN-centered subset is
\[
  \mathcal{D}_{\mathrm{BN}}=\{r_i:\{B,N\}\subseteq\mathrm{elements}(c_i)\ \mathrm{or}\ c_i\in\mathcal{C}_{\mathrm{BN}}\}.
\]
Deduplication, property alignment, dimensionality labels, and source-version tracking will be implemented before modelling.

\subsection{WP2: Formula-Only and Structure-Resolved Modelling}
The modelling layer will compare descriptor baselines, composition neural models, and structure-resolved graph models using established tools such as pymatgen, CGCNN, MEGNet, ALIGNN, CrabNet, and CHGNet~\cite{ong2013_pymatgen,xie2018_cgcnn,chen2019_megnet,choudhary2021_alignn,wang2021_crabnet,deng2023_chgnet}. The application vector is
\[
  \mathbf{g}(c,s)=
  [\gapscore,\;-\ehull,\;\rho_{\mathrm{dielectric}},\;\rho_{\mathrm{UV}},
  \rho_{\mathrm{2D}},\;-\unc].
\]
The band-gap term will be treated as a target-window score rather than a monotonic objective:
\[
  \gapscore(c)=
  \exp\!\left[-\frac{(\hat{\eg}(c)-E^*_{\mathrm{g,app}})^2}{2\sigma_{\mathrm{app}}^2}\right]
  \cdot\eta_{\mathrm{direct}}(c).
\]
The scores $\rho_{\mathrm{dielectric}}$, $\rho_{\mathrm{UV}}$, and $\rho_{\mathrm{2D}}$ are proxy scores derived from computed properties, dimensionality metadata, and rule-based target criteria. Predictive uncertainty will be estimated by ensembles and calibrated against validation errors:
\[
  \hat{\unc}_{\mathrm{cal}}(c)=
  \alpha\,\mathrm{Std}_{m\in\mathcal{M}}\!\left[f_m(c)\right]+
  \beta\,\widehat{\mathrm{err}}_{\mathrm{val}}(c).
\]
If structure-resolved models do not outperform composition or descriptor baselines under BN-family-aware splits, graph-model outputs will be used as exploratory evidence rather than as the primary ranking signal.

\subsection{WP3: Candidate Prioritization and Exploration Protocol}
Candidate decisions will use both predictive value and reliability. Domain distance will be measured from composition descriptors, structural descriptors, and learned embeddings:
\[
  \domain(c)=\min_{r_i\in\mathcal{D}_{\mathrm{BN}}}
  \left\|\phi(c)-\phi(c_i,s_i)\right\|_{\Sigma}.
\]
Novelty will be estimated by distance from known BN-family compounds, while application relevance will be assigned by matching the candidate to wide-gap, dielectric, or 2D-support criteria. The action label is
\[
  a(c)\in\{\mathrm{control},\mathrm{priority},\mathrm{explore},\mathrm{hold}\},
\]
where priority candidates proceed to structure handoff, exploratory candidates are retained as high-risk families, controls provide reference checks, and hold candidates are deferred until additional evidence is available. The operational protocol is:
\[
  \mathrm{pool}\rightarrow\mathrm{property\ prediction}\rightarrow
  \mathrm{uncertainty/domain\ estimation}\rightarrow
  \mathrm{filtering}\rightarrow\mathrm{action\ label}.
\]
This design makes the project a usable AI decision pipeline rather than a list of disconnected materials-ML tasks.

\subsection{WP4: Structure Handoff and Validation}
Selected formulas will first be mapped to candidate structures by prototype reuse, structure transfer, and controlled substitution. Generative models such as CDVAE will be considered only when prototype-based routes are insufficient~\cite{xie2022_cdvae}. Validation will include charge-neutrality checks, structural sanity checks, geometry relaxation where feasible, energy-above-hull re-estimation,
\[
  \ehull(c,s)=E(c,s)-E_{\mathrm{hull}}(\mathrm{composition}(c)),
\]
and comparison with known BN-family references. Stability will be handled separately for bulk-like and 2D candidates: bulk-like entries will use formation-energy or energy-above-hull proxies, while 2D entries will additionally use dimensionality, database provenance, and available 2D stability indicators. Evaluation will use random, grouped-by-formula, and BN-family-aware splits, with MAE, RMSE, $R^2$, calibration error, Kendall $\tau$, Spearman $\rho$, and top-$k$ stability as metrics.

\section{Significance of the Proposed Research and Major Deliverables}
The significance of this project is not merely applying existing models to BN. Its contribution is a BN-specific AI exploration framework: a curated data layer, family-aware evaluation, calibrated uncertainty, application-aware ranking, and prototype-first structure handoff. This is directly relevant to BN research in ultraviolet materials, dielectric layers, 2D heterostructures, and related functional-material discovery~\cite{watanabe2004_hbn_uv,cassabois2016_hbn_indirect,dean2010_hbn_graphene}. It also produces a reusable AI-for-materials workflow that can be extended to other focused materials families.

\textbf{Major deliverables:}
\begin{enumerate}[leftmargin=*]
  \item \textbf{D1: BN-centered dataset} with formula, structure, target-property, dimensionality, and provenance fields from at least three public sources.
  \item \textbf{D2: Benchmark report} comparing descriptor, composition, and structure-resolved models under random, grouped-by-formula, and BN-family-aware splits.
  \item \textbf{D3: Ranked candidate table} of at least 30 BN-related candidate families, including predicted band gap, stability proxy, uncertainty, domain distance, novelty score, application label, and recommended action.
  \item \textbf{D4: Structure hypothesis package} for 5--10 priority candidates, including CIF/POSCAR files, provenance records, predicted properties, uncertainty scores, and validation notes.
  \item \textbf{D5: Research outputs} including technical report, figures, demo interface, and manuscript-ready analysis for PI and collaborator review.
\end{enumerate}

\section{Ethics and Safety}
The initial project is computational and data-driven. It uses public databases, open-source software, models, and generated candidate structures. No human subjects, living animals, wet-lab synthesis, biological samples, or radiation work will be conducted in the initial computational phase. If later experimental synthesis or regulated laboratory procedures are introduced, the responsible experimental party will obtain the required safety approvals before that work begins.

\section{Selected References}
\small
\begin{thebibliography}{99}
\setlength{\itemsep}{-1pt}
\bibitem{zhou2019_2dmatpedia} J. Zhou et al., ``2DMatPedia: an open computational database of two-dimensional materials from top-down and bottom-up approaches,'' \textit{Scientific Data}, 6, 86, 2019.
\bibitem{choudhary2020_jarvis} K. Choudhary et al., ``JARVIS-Tools: open-access software for atomistic data-driven materials design,'' \textit{npj Computational Materials}, 6, 173, 2020.
\bibitem{dunn2020_matbench} A. Dunn et al., ``Benchmarking materials property prediction methods: the Matbench test set and Automatminer reference algorithm,'' \textit{npj Computational Materials}, 6, 138, 2020.
\bibitem{ward2018_matminer} L. Ward et al., ``Matminer: An open source toolkit for materials data mining,'' \textit{Computational Materials Science}, 152, 60--69, 2018.
\bibitem{jain2013_mp} A. Jain et al., ``The Materials Project: a materials genome approach to accelerating materials innovation,'' \textit{APL Materials}, 1, 011002, 2013.
\bibitem{haastrup2018_c2db} S. Haastrup et al., ``The Computational 2D Materials Database,'' \textit{2D Materials}, 5, 042002, 2018.
\bibitem{ong2013_pymatgen} S. P. Ong et al., ``Python Materials Genomics (pymatgen),'' \textit{Computational Materials Science}, 68, 314--319, 2013.
\bibitem{xie2018_cgcnn} T. Xie and J. C. Grossman, ``Crystal graph convolutional neural networks for materials property prediction,'' \textit{Physical Review Letters}, 120, 145301, 2018.
\bibitem{chen2019_megnet} C. Chen et al., ``Graph networks as a universal machine learning framework for molecules and crystals,'' \textit{Chemistry of Materials}, 31, 3564--3572, 2019.
\bibitem{choudhary2021_alignn} K. Choudhary and B. DeCost, ``Atomistic line graph neural network for improved materials property predictions,'' \textit{npj Computational Materials}, 7, 185, 2021.
\bibitem{wang2021_crabnet} A. Y.-T. Wang et al., ``CrabNet for explainable deep learning in materials science,'' \textit{npj Computational Materials}, 7, 77, 2021.
\bibitem{deng2023_chgnet} B. Deng et al., ``CHGNet as a pretrained universal neural network potential,'' \textit{Nature Machine Intelligence}, 5, 1031--1041, 2023.
\bibitem{xie2022_cdvae} T. Xie et al., ``Crystal diffusion variational autoencoder for periodic material generation,'' \textit{ICLR}, 2022.
\bibitem{watanabe2004_hbn_uv} K. Watanabe, T. Taniguchi and H. Kanda, ``Direct-bandgap properties and ultraviolet lasing of hexagonal boron nitride,'' \textit{Nature Materials}, 3, 404--409, 2004.
\bibitem{cassabois2016_hbn_indirect} G. Cassabois, P. Valvin and B. Gil, ``Hexagonal boron nitride is an indirect bandgap semiconductor,'' \textit{Nature Photonics}, 10, 262--266, 2016.
\bibitem{dean2010_hbn_graphene} C. R. Dean et al., ``Boron nitride substrates for high-quality graphene electronics,'' \textit{Nature Nanotechnology}, 5, 722--726, 2010.
\end{thebibliography}
\normalsize

\section*{Project Details}
\textbf{Project Title:} AI-Powered Boron Nitride Material Exploration

\textbf{PI (Dept):} Prof. Chen MA (Department of Computer Science)

\textbf{Budget Breakdown}
\begin{center}
\small
\begin{tabular}{>{\raggedright\arraybackslash}p{0.19\linewidth}>{\raggedright\arraybackslash}p{0.55\linewidth}>{\raggedright\arraybackslash}p{0.16\linewidth}}
\toprule
\textbf{Budget Item} & \textbf{Detail breakdown and justification} & \textbf{Amount (HKD)}\\
\midrule
Staffing & Student helper / research assistant support for data cleaning, benchmark scripts, candidate-table verification, figure generation, documentation, and demo maintenance. & \blank\\
Equipment & GPU cloud hours for graph-model training and ensemble inference; storage for downloaded structures, generated CIF/POSCAR files, benchmark logs, and model checkpoints. & \blank\\
General expenses & Cloud credits, repository backup, software services, lightweight demo hosting, and data-management costs. & \blank\\
Conference & Registration and presentation materials for AI-for-science, materials informatics, or applied computing venues. & \blank\\
Overseas Travel & Optional PI-approved collaborator visit, technical meeting, or conference travel for candidate-review discussion. & \blank\\
\midrule
 & \textbf{Total Funding} & \blank\\
\bottomrule
\end{tabular}
\end{center}

\textbf{Collaborator(s):} Shenzhen Zhongneng Tianyi Technology Co., Ltd.

\end{document}
